\documentclass[11pt]{article}

% -----------------------------
% Packages
% -----------------------------
\usepackage[utf8]{inputenc}
\usepackage[T1]{fontenc}
\usepackage{lmodern}

\usepackage{geometry}
\geometry{a4paper, margin=1in}

\usepackage{setspace}
\onehalfspacing

\usepackage{amsmath, amssymb}
\usepackage{booktabs}
\usepackage{graphicx}

\usepackage{natbib}
\bibliographystyle{apalike}

\usepackage{hyperref}
\hypersetup{
    colorlinks=true,
    citecolor=blue,
    linkcolor=blue,
    urlcolor=blue
}

% -----------------------------
% Title
% -----------------------------
\title{Systematic and Idiosyncratic Components of Relative Signed Jump Variation}
\author{}
\date{}

\begin{document}

\maketitle

\section{Methodology}

\subsection{Decomposition of RSJ into Systematic and Idiosyncratic Components}

To distinguish between common market-wide asymmetric volatility
conditions and firm-specific asymmetry, we decompose the Relative
Signed Jump (RSJ) measure into systematic and idiosyncratic components
using a factor-model framework.

\subsection{Construction of Firm-Level RSJ}

Following \citet{bollerslev2019}, realized volatility is decomposed into
positive and negative semivariance components using intraday returns.
Let $r_{i,t,j}$ denote the $j$-th intraday return for stock $i$ on day $t$.

Daily realized variance is defined as

\[
RV_{i,t} = \sum_{j=1}^{n} r_{i,t,j}^{2}.
\]

Positive and negative realized semivariance are defined as

\[
RV^+_{i,t} =
\sum_{j=1}^{n}
r_{i,t,j}^{2} \mathbf{1}(r_{i,t,j} > 0),
\]

\[
RV^-_{i,t} =
\sum_{j=1}^{n}
r_{i,t,j}^{2} \mathbf{1}(r_{i,t,j} < 0).
\]

The Relative Signed Jump measure is then defined as

\[
RSJ_{i,t} =
\frac{RV^+_{i,t} - RV^-_{i,t}}{RV_{i,t}}.
\]

This measure captures the relative contribution of positive and
negative intraday return variation to total realized volatility.

\subsection{Construction of an Aggregate RSJ Factor}

To capture common market-wide asymmetry in volatility, we construct an
aggregate RSJ factor defined as the cross-sectional weighted average
of firm-level RSJ measures:

\[
RSJ^M_t =
\sum_{i=1}^{N_t}
w_{i,t} RSJ_{i,t},
\]

where $w_{i,t}$ denotes the weight of stock $i$ in period $t$.
In the baseline specification, these weights correspond to
market-capitalization weights, although equal-weighted aggregates will
also be considered as a robustness check.

This aggregate measure captures the common component of asymmetric
volatility across firms.

\subsection{Factor-Model Decomposition}

Firm-level RSJ is decomposed into systematic and idiosyncratic
components using the following factor model:

\[
RSJ_{i,t} =
\alpha_i +
\beta_i RSJ^M_t +
\varepsilon_{i,t}.
\]

The fitted value represents the systematic component:

\[
RSJ^{sys}_{i,t} =
\beta_i RSJ^M_t,
\]

while the residual represents the idiosyncratic component:

\[
RSJ^{idio}_{i,t} =
\varepsilon_{i,t}.
\]

The systematic component captures the part of firm-level asymmetric
volatility related to common market conditions, whereas the
idiosyncratic component captures firm-specific asymmetry.

\subsection{Return Predictability Tests}

To examine which component drives return predictability, we estimate
cross-sectional regressions of the form

\[
R_{i,t+1} =
\alpha_t +
\gamma_{1,t} RSJ^{sys}_{i,t} +
\gamma_{2,t} RSJ^{idio}_{i,t} +
\delta_t' X_{i,t} +
\varepsilon_{i,t+1},
\]

where $R_{i,t+1}$ denotes the future return of stock $i$ and $X_{i,t}$
is a vector of control variables commonly used in the cross-sectional
asset pricing literature, such as size, book-to-market, momentum, and
idiosyncratic volatility.

If the predictive power of RSJ is mainly driven by the systematic
component, this would support a risk-based interpretation in which
investors require compensation for exposure to common downside
asymmetry conditions. Conversely, if the idiosyncratic component
remains predictive, the evidence would be more consistent with a
stock-specific mispricing interpretation.

\bibliography{references}

\end{document}